\subsection{FUNDAMENTAL SYSTEMS OF NEIGHBOURHOODS; BASES OF A TOPOLOGY}

DEFINITION 5. In a topological space X, a fundamental system of neighbourhoods of a point x (resp. of a subset A of X) is any set S of neighbourhoods of x (resp. A) such that for each neighbourhood V of x (resp. A) there is a neighbourhood W ∈ S such that W ⊂ V.

If \(\mathfrak{S}\) is a fundamental system of neighbourhoods of a subset \(A\) of \(X\) , then every finite intersection of sets of \(\mathfrak{S}\) contains a set of \(\mathfrak{S}\) .

Examples. 1) In a discrete space (no. 1) the set \(\{x\}\) alone constitutes a fundamental system of neighbourhoods of the point \(x\) .

\begin{enumerate}
\def\labelenumi{\arabic{enumi})}
\setcounter{enumi}{1}
\tightlist
\item
  On the rational line Q the set of all open intervals containing a point x is a fundamental system of neighbourhoods of this point. So is the set of open intervals \(]x - 1/n, x + 1/n[\) , and the set of closed intervals \([x - 1/n, x + 1/n]\) , where n runs through all the integers \textgreater0, or through any infinite strictly increasing sequence of integers \textgreater0. * There are analogous results for the real line. *
\end{enumerate}

DEFINITION 6. A base of the topology of a topological space X is any set B of open subsets of X such that every open subset of X is the union of sets belonging to B.

PROPOSITION 3. If X is a topological space, then for a set B of open subsets of X to be a base of the topology of X it is necessary and sufficient that for each \(x \in X\) the set of \(V \in B\) such that \(x \in V\) is a fundamental system of neighbourhoods of x.

It is clear that the condition is necessary. Conversely, if it is satisfied, then given any open set U and any \(x \in U\) there is an open set \(V_{x} \in B\) such that \(x \in V_{x} \subset U\) . The union of the sets \(V_{x}\) for \(x \in U\) is therefore equal to U. This completes the proof.

Examples. 1) The discrete topology has as a base the set of subsets of X which consist of a single point.

\begin{enumerate}
\def\labelenumi{\arabic{enumi})}
\setcounter{enumi}{1}
\tightlist
\item
  The set of bounded open intervals is by definition a base of the topology of the rational line (no. 2). * Likewise, the set of bounded open intervals is a base of the topology of the real line. *
\end{enumerate}

